% !TEX program = xelatex
\documentclass[10pt,letterpaper]{article}
\usepackage{../ac-paper-gear-guess}
% Purchase-link layer for aesthetic-computer-media.
%
% The checked-in defaults are direct, non-affiliate manufacturer links.
% After Amazon Associates approval:
%   1. replace each value with the exact SiteStripe short link for that SKU;
%   2. change \gearaffiliatefalse to \gearaffiliatetrue;
%   3. rebuild and click-test the PDF.
%
% Do not invent a tag or use another publisher's Associates ID.

\gearaffiliatefalse

\newcommand{\buyOWCExpressOneMTwo}{https://www.owc.com/solutions/express-1m2}
\newcommand{\buySamsungNineNineZeroPro}{https://www.samsung.com/us/memory-storage/nvme-ssd/990-pro-pcie-4-0-nvme-ssd-1tb-sku-mz-v9p2t0b-am/}
\newcommand{\buySanDiskProGForty}{https://www.sandisk.com/products/ssd/external-ssd/sandisk-professional-pro-g40-usb-3-2-ssd}
\newcommand{\buySamsungTNine}{https://www.samsung.com/us/memory-storage/portable-ssd/portable-ssd-t9-usb-3-2-4tb-black-sku-mu-pg4t0b-am/}


\gearsetshorttitle{aesthetic-computer-media}
\hypersetup{
  pdftitle={aesthetic-computer-media — a gear guess},
  pdfsubject={A portable high-speed media workspace for the Aesthetic Computer fleet},
  pdfkeywords={external SSD, USB4, media workflow, Aesthetic Computer}
}

\begin{document}

\gearGuessTitle
  {aesthetic-computer-media}
  {one fast portable workspace for neo, blueberry, and whatever is next}
  {22 July 2026}
  {guess 1}

\begin{bestguess}
Buy the \textbf{OWC Express 1M2 USB4 enclosure} and put a
\textbf{Samsung 990 PRO 2TB NVMe} inside it. Give the physical object the long
name \texttt{aesthetic-computer-media}; format its shared volume as
\texttt{AC\_MEDIA}. This is the useful middle: near-internal-disk speed on a
USB4/Thunderbolt host, ordinary USB-C fallback, a replaceable SSD, and no wall
wart. Start at 2TB because the drive is the \emph{active table}, not the archive.
DigitalOcean Spaces remains the shelf.
\end{bestguess}

\subsection{canonical object plate}
\gearproduct
  {figures/owc-express-1m2.png}
  {OWC Express 1M2}
  {preferred enclosure}
  {\$119.99 observed}
  {https://www.owc.com/solutions/express-1m2}
  {\buyOWCExpressOneMTwo}
\hfill
\gearproduct
  {figures/samsung-990-pro.jpg}
  {Samsung 990 PRO 2TB}
  {preferred blade}
  {\$390 allowance}
  {https://www.samsung.com/us/memory-storage/nvme-ssd/990-pro-pcie-4-0-nvme-ssd-1tb-sku-mz-v9p2t0b-am/}
  {\buySamsungNineNineZeroPro}

\vspace{0.12in}
\gearproduct
  {figures/sandisk-pro-g40.png}
  {SanDisk Professional PRO-G40}
  {rugged alternate}
  {\$919.99 observed}
  {https://www.sandisk.com/products/ssd/external-ssd/sandisk-professional-pro-g40-usb-3-2-ssd}
  {\buySanDiskProGForty}
\hfill
\gearproduct
  {figures/samsung-t9.png}
  {Samsung Portable SSD T9}
  {one-piece alternate}
  {\$799.99 observed}
  {https://www.samsung.com/us/memory-storage/portable-ssd/portable-ssd-t9-usb-3-2-4tb-black-sku-mu-pg4t0b-am/}
  {\buySamsungTNine}

\gearaffiliatedisclosure

\subsection{capacity rail}

\begin{tabularx}{\linewidth}{@{}Y R R@{}}
\toprule
\textbf{decision rail} & \textbf{2TB now} & \textbf{4TB step-up} \\
\midrule
Enclosure allowance & \$120 & \$120 \\
NVMe allowance & \$390 & \$700--750 \\
Working total & \textbf{\$510} & \textbf{\$820--870} \\
\bottomrule
\end{tabularx}
{\footnotesize\color{acgray}Prices are observed/working allowances checked 22 July 2026, before tax and shipping. They are not quotes. The 4TB number is a range because flash pricing moves quickly.}

\section{buy this}

\begin{gearcard}{1 / the enclosure}
\gearprice{OWC Express 1M2 — USB4, 40Gb/s}{\$119.99 observed}
Bus-powered aluminum enclosure; a short certified 40Gb/s cable is included.
OWC reports roughly 3,189 MB/s on Apple silicon and up to 3,836 MB/s on
USB4/Thunderbolt 4 PCs, with roughly 990 MB/s fallback on a 10Gb/s USB-C link.
The enclosure is the durable purchase. The blade inside can change later.
\end{gearcard}

\begin{gearcard}{2 / the blade}
\gearprice{Samsung 990 PRO NVMe — 2TB}{\$390 allowance}
The bare drive is rated up to 7,450 MB/s read and 6,900 MB/s write, so the USB4
enclosure---not the flash---sets portable speed. That headroom is still useful:
the blade is never the obvious bottleneck, and it can move into a workstation
or a future enclosure. Samsung gives the line a five-year limited warranty;
the 4TB part is rated for 2,400 TBW.
\end{gearcard}

\begin{gearwarning}[THE PRICE RULE]
Do not pay extra for the 990 PRO name if another reputable TLC NVMe with a
five-year warranty makes the 2TB build materially cheaper. The enclosure caps
the result near 4 GB/s. The exact blade is negotiable; USB4, thermals, and a
replaceable drive are not. Step up to 4TB only when the live working set is
regularly above 1TB or the price delta falls below roughly \$250.
\end{gearwarning}

\section{the other two}

\begin{gearcard}{rugged, sealed, simple / SanDisk Professional PRO-G40 4TB}
\gearprice{Best when the drive will travel hard}{\$919.99 observed}
Thunderbolt 3 performance is advertised up to 3,000 MB/s read and 2,500 MB/s
write. The enclosure is IP68-rated, drop-rated to 3m, crush-resistant to
4,000lb, and carries a five-year warranty. Pick this instead when weather,
bags, stages, and other people's hands matter more than replaceable internals.
The official compatibility list names macOS and Windows; test the intended
Fedora machine before making it the fleet standard.
\end{gearcard}

\begin{gearcard}{one-piece fallback / Samsung T9 4TB}
\gearprice{Best only at a real discount}{\$799.99 observed}
Small, light, and broadly available. Its 2,000 MB/s headline needs a USB 3.2
Gen 2x2 host. On a 10Gb/s USB link, expect the neighborhood of 1 GB/s instead.
An independent sustained-write test also found roughly 900 MB/s after the
cache emptied. Pick it if it falls well below the modular build and maximum
Mac speed is not the point.
\end{gearcard}

\begin{gearwarning}[SKIP FOR THIS FLEET]
Skip the newer 80Gb/s USB4 v2 / Thunderbolt 5 enclosure for now. Its advantage
appears only on hosts with those ports. Mixed Macs and PCs make the 40Gb/s 1M2
the more portable kind of fast.
\end{gearwarning}

\section{port reality}

\begin{tabularx}{\linewidth}{@{}p{0.27\linewidth}Y p{0.22\linewidth}@{}}
\toprule
\textbf{host link} & \textbf{what it means} & \textbf{working ceiling} \\
\midrule
USB4 / Thunderbolt 4 & The recommended build gets its full reason to exist. & about 3.2--3.8 GB/s \\
Thunderbolt 3 & Fast path; controller and host still matter. & about 2.5--3.0 GB/s \\
USB-C, 10Gb/s & Compatible fallback, not failure. & about 1.0 GB/s \\
USB 3.2 Gen 2x2 & Helps the T9 only when the exact host supports it. & up to 2.0 GB/s \\
\bottomrule
\end{tabularx}

The connector shape proves nothing. Before blaming the disk, identify the host
port, use the included short cable, and run one large sequential read/write
test. Keep the aluminum enclosure exposed during long writes; it is the heat
sink.

\section{day-one score}

\begin{gearcard}{label + format}
\gearcheck Put \texttt{aesthetic-computer-media} on the physical case.\par
\gearcheck Initialize a GUID partition table and one exFAT volume named
\texttt{AC\_MEDIA}. exFAT's volume label is limited to eleven Unicode
characters; the short name is deliberate.\par
\gearcheck Leave hardware-password features off if Fedora must read the drive.
Encrypt sensitive project bundles separately instead of making the whole
working disk depend on vendor unlock software.\par
\gearcheck Always eject. exFAT is the shared language here, not a journaled
filesystem and not a backup.
\end{gearcard}

\begin{gearcard}{same object, predictable places}
\begin{tabularx}{\linewidth}{@{}p{0.24\linewidth}Y@{}}
macOS & \texttt{/Volumes/AC\_MEDIA} \\
Windows & reserve \texttt{M:\textbackslash} \\
Fedora & \texttt{/run/media/jas/AC\_MEDIA} \\
\end{tabularx}
\vspace{0.25em}
Create only four roots at first:
\texttt{active/}, \texttt{inbox/<machine>/}, \texttt{exports/}, and
\texttt{manifests/}. Projects move between machines; machine names belong only
in \texttt{inbox/}.
\end{gearcard}

\begin{gearcard}{the shelf loop}
\begin{enumerate}
  \item Copy new media to \texttt{inbox/<machine>/}; keep the source intact.
  \item Promote the working set into \texttt{active/<project>/}.
  \item Write a SHA-256 manifest under \texttt{manifests/<project>/}.
  \item \texttt{rclone copy} the project to the private Spaces
        \texttt{media/} prefix.
  \item Run \texttt{rclone check --download}. Delete the ingest source only
        after the check succeeds.
  \item Enable Spaces versioning before treating cloud-only media as recoverable.
\end{enumerate}
\end{gearcard}

\begin{bestguess}[WHAT WOULD CHANGE THE GUESS]
A fleet-wide move to Thunderbolt 5 would justify an 80Gb/s enclosure. A
working set consistently above 1TB would justify 4TB. Regular field abuse would
move the PRO-G40 into first place. A strong sale on a reputable TLC blade would
change the SSD model without changing the system. Until one of those happens,
the 1M2 + 2TB NVMe is the clean answer.
\end{bestguess}

\section{source shelf}

Every product number above is a dated observation, not a promise from the
future. Product capabilities come from the manufacturer unless an independent
measurement is named.

\begin{multicols}{2}
\raggedcolumns
\gearcompactsource{OWC Express 1M2 product page}{checked 22 July 2026}{https://www.owc.com/solutions/express-1m2}
\gearcompactsource{Samsung 990 PRO product page}{checked 22 July 2026}{https://www.samsung.com/us/memory-storage/nvme-ssd/990-pro-pcie-4-0-nvme-ssd-1tb-sku-mz-v9p2t0b-am/}
\gearcompactsource{Samsung 990 PRO data sheet, revision 2.0}{checked 22 July 2026}{https://download.semiconductor.samsung.com/resources/data-sheet/samsung_nvme_ssd_990_pro_datasheet_rev.2.0.pdf}
\gearcompactsource{SanDisk Professional PRO-G40 product page}{checked 22 July 2026}{https://www.sandisk.com/products/ssd/external-ssd/sandisk-professional-pro-g40-usb-3-2-ssd}
\gearcompactsource{Samsung Portable SSD T9 product page}{checked 22 July 2026}{https://www.samsung.com/us/memory-storage/portable-ssd/portable-ssd-t9-usb-3-2-4tb-black-sku-mu-pg4t0b-am/}
\gearcompactsource{Best Buy — Samsung T9 4TB listing}{checked 22 July 2026}{https://www.bestbuy.com/product/samsung-t9-portable-ssd-4tb-up-to-2000mb-s-usb-3-2-gen2-black/6559268}
\gearcompactsource{PC Gamer — Samsung T9 sustained-write review}{checked 22 July 2026}{https://www.pcgamer.com/hardware/ssds/samsung-t9-external-ssd-review/}
\gearcompactsource{Microsoft exFAT specification}{checked 22 July 2026}{https://learn.microsoft.com/en-us/windows/win32/fileio/exfat-specification}
\gearcompactsource{Microsoft filesystem functionality comparison}{checked 22 July 2026}{https://learn.microsoft.com/en-us/windows/win32/fileio/filesystem-functionality-comparison}
\gearcompactsource{DigitalOcean — enable Spaces versioning}{checked 22 July 2026}{https://docs.digitalocean.com/products/spaces/how-to/enable-versioning/}
\gearcompactsource{DigitalOcean — transfer and verify with rclone}{checked 22 July 2026}{https://docs.digitalocean.com/products/spaces/how-to/transfer-between-regions/}
\gearcompactsource{Canonical product image ledger}{checked 22 July 2026}{https://github.com/whistlegraph/aesthetic-computer/tree/main/papers/gear-guess-aesthetic-computer-media/figures/product-images.md}
\end{multicols}

\vfill
{\color{acpink}\rule{\linewidth}{1pt}}\par
{\small\sffamily
\textbf{Gear Guess 1 / aesthetic-computer-media.}
This paper names a working decision. Recheck price, stock, host ports, and
warranty terms at purchase time. The operating idea is more stable than the
specific SKU: fast replaceable local table; verified private cloud shelf.}

\end{document}
